\documentclass[12pt]{article}
\usepackage[utf8]{inputenc}
\usepackage[T1]{fontenc}
\usepackage{hyperref}
\title{Zizilivakan Language Resources for K-12 Curriculum Development}
\author{Compiled by Claude Sonnet 5.5 (claude-sonnet-5-5) for Andr\'as Kornai}
\date{}
\begin{document}
\maketitle

\section{Introduction}
Zizilivakan (ISO 639-3: ziz) is a Chadic language, within the Afro-Asiatic family, spoken in Adamawa State, Nigeria, and the Far North Region of Cameroon. Other names include Ziziliveken, Ziliva and Fali of Jilbu. Speaker counts differ by source: Wikipedia gives about 6,000, Joshua Project lists about 9,000 in Nigeria and about 300 in Cameroon, and the Endangered Languages Project gives 2,800. UNESCO rates it as definitely endangered. Only general reference pages turned up in the searches.

\section{Primary Educational Resources}
No primers, textbooks or school programmes for Zizilivakan turned up in the searches.

\section{Linguistic and Reference Materials}
\begin{itemize}
\item Wikipedia, Zizilivakan language: \url{https://en.wikipedia.org/wiki/Zizilivakan_language}
\item UNESCO Atlas of the World's Languages, Zizilivakan (Cameroon): \url{https://en.wal.unesco.org/countries/cameroon/languages/zizilivakan}
\item Endangered Languages Project entry: \url{https://www.endangeredlanguages.com/node/89330}
\end{itemize}

\section{Religious and Community Resources}
\begin{itemize}
\item Joshua Project, Fali of Jibu in Nigeria. It describes most schooling as in English or Hausa: \url{https://www.joshuaproject.net/index.php/people_groups/11182/NI}
\item Joshua Project, Fali of Jibu in Cameroon: \url{https://www.joshuaproject.net/index.php/people_groups/11182/CM}
\end{itemize}

\section{Digital and Audio Resources}
No audio, video or software for Zizilivakan turned up in the searches.

\section{Curriculum Development Context}
Little published material exists for Zizilivakan. A curriculum would need new documentation, an orthography and teaching materials, ideally developed with the community. English and Hausa are the current languages of schooling in Nigeria.
\end{document}